\SourceLeaf{232}%
we choose a definite side on which we imagine the figure to~lie. On the auxiliary sphere we shall always use the exterior side, facing away from the centre. On the curved surface we may likewise use the exterior side or the side regarded as exterior, or more precisely the side from which we imagine the normal rising. Clearly the similarity of the figures is unaffected if both the figure on the curved surface and the normal are transferred to the opposite side, provided that the image of the figure is always drawn on the same side of the sphere.

We also extend to the integral curvature of a finite figure on the curved surface the positive or negative sign that we assign to the \textit{measure} of curvature according to the position of an infinitely small figure. But if we undertake to treat the subject in full generality, a few points need clarification, which we shall touch on only briefly here. As long as the figure on the curved surface is such that each point inside it corresponds to a \textit{distinct} point on the sphere, the definition needs no further explanation. When this condition fails, we shall have to count certain parts of the figure on the sphere twice or more, so that they may add up or cancel according to whether their position is similar or reversed. In such a case the simplest way is to imagine the figure on the curved surface divided into parts that each satisfy the condition when considered separately. We assign each part its integral curvature, with the magnitude determined by the area of the corresponding figure on the sphere and the sign by its position. Finally we assign the whole figure the integral curvature obtained by adding the integral curvatures of its individual parts. In general, then, the integral curvature of a figure is \(=\int k\,\mathrm{d}\sigma\), where \(\mathrm{d}\sigma\) denotes the element of area of the figure and \(k\) the measure of curvature at any point. The main points concerning the geometric representation of this integral are as follows. The boundary of the figure on the curved surface, subject to the restriction of Article 3, will always correspond to a line on the sphere that returns to itself. If this line nowhere intersects itself, it will divide the whole sphere into two parts, one of which corresponds to the figure on the curved surface. The area of that part gives the integral curvature of the figure on the curved surface when taken positively or negatively according to whether that part lies relative to its own boundary in the same way as the original figure relative to its boundary or in the reverse way. But whenever the line crosses itself once or more, it gives a complicated figure, to which we may nevertheless assign a definite area just as legitimately as to figures without nodes. Properly understood, this area will always give the correct value of the integral curvature. We must, however, %
